% template.tex --- starting point for a set of lecture notes in
% lecturenotes.cls. Every element of the format appears once, filled with
% placeholder text: keep the skeleton, replace the content. The rules behind
% each element are in docs/formatting-guide.md.
%
% Build: pdflatex template.tex (run twice).
\documentclass{lecturenotes}

\title{Title of the Note}
\subtitle{The one question this note answers?}
\author{Author Name}
\series{Series name --- lecture notes}
% \runningtitle{Short title}  % left-hand header, if the title is long
% \date{...}                  % defaults to the current month and year

\begin{document}

\maketitle
\tableofcontents

% ===========================================================================
% Front matter: one-line description, three lead paragraphs, reading map.
% ===========================================================================
One sentence stating the subject and that it is developed from first
principles.

\paragraph{Prerequisites.} The mathematics the reader is assumed to have,
listed concretely. \textbf{Nothing beyond that is assumed.} Name any
advanced machinery the notes develop for themselves rather than assume.

\paragraph{What you should be able to do afterwards.} Three or four concrete
capabilities, phrased as verbs: derive this result rather than recall it,
state exactly what that theorem does and does not assert, apply the method to
the case study.

\paragraph{Companion notes.} How this note relates to its siblings in the
series, and which sections carry the connections (for example
\cref{sec:second-sub}).

\separator

\subsubsection*{Reading map}

\paragraph{Part 0 --- Vocabulary.} \Cref{sec:everything} fixes notation and
states every result used but not derived. Skip it and refer back as needed.

\paragraph{Part I --- Foundations.} \Cref{sec:first} sets up the object of
study; \cref{sec:second} derives the main result.

\paragraph{Part II --- Practice.} \Cref{sec:case} is the worked case;
\cref{sec:limits} states the limits. \Cref{app:checks} records the
numerical checks.

% ===========================================================================
\part{Notation and vocabulary}
% ===========================================================================

\section{Everything used but not derived}\label{sec:everything}

These notes are self-contained. \Cref{sec:notation} fixes notation;
\cref{sec:results} states the results used, with a proof where one is short
and an explicit statement of trust where it is not. Nothing later uses a
symbol, result or term that does not appear here or at its point of first
use.

\subsection{Notation}\label{sec:notation}

\begin{tabularx}{\linewidth}{@{}l L@{}}
\toprule
Symbol & Meaning \\
\midrule
$x$, $y$ & placeholder: the primary variables \\ \rowrule
$\theta$ & placeholder: the parameter being estimated \\ \rowrule
$f(x;\theta)$ & placeholder: the model, with a cross-reference to where it is defined (\cref{def:first}) \\ \rowrule
$\mathbb E$, $\mathrm{Var}$ & expectation, variance \\ \rowrule
$\log$ & natural logarithm throughout \\ \rowrule
$\blacksquare$ & end of proof \\
\bottomrule
\end{tabularx}

\subsection{Mathematical results used}\label{sec:results}

\begin{definition}[name of the concept]\label{def:first}
Placeholder definition. The term being defined is set in
\textbf{bold} at its point of definition; the rest of the statement is
upright text with inline mathematics such as $f(x;\theta)$.
\end{definition}

\begin{lemma}[name of the lemma]\label{lem:first}
Placeholder statement of a short auxiliary result, e.g.\ if $a \le b$ then
$\varphi(a) \le \varphi(b)$ for every non-decreasing $\varphi$.
\end{lemma}

\begin{proof}
Placeholder proof, ending with the filled square the class supplies
automatically.
\end{proof}

\begin{fact}[name of the fact]\label{fact:first}
Placeholder statement of a standard result used without proof. Say where a
proof can be found~\cite{ref-book}; a Fact is the only statement kind
permitted to appear without one.
\end{fact}

% ===========================================================================
\part{Foundations}
% ===========================================================================

\section{First topic}\label{sec:first}

Placeholder paragraph motivating the section: what problem it solves and why
the obvious approach fails. Display equations pad their top-level relations
with thick spaces:
\[
  y \;=\; f(x;\theta) \;+\; \varepsilon,
  \qquad \mathbb E[\varepsilon] \;=\; 0.
\]

\begin{proposition}[name of the proposition]\label{prop:first}
Placeholder statement of an intermediate result.
\end{proposition}

\begin{proof}
Placeholder proof, citing \cref{lem:first} and \cref{fact:first} by name
rather than by bare number.
\end{proof}

\begin{summary}
Placeholder plain-language restatement: two or three sentences saying what
the derivation above established, without symbols where possible.
\end{summary}

\section{Second topic}\label{sec:second}

\subsection{First subsection}

\begin{theorem}[name of the main result, Author~\cite{ref-article}]\label{thm:main}
Placeholder statement of the main theorem. An equation that is referred to
later is numbered and labelled:
\begin{equation}\label{eq:main}
  F(\theta) \;=\; \mathbb E\big[\thinspace\ell(y, f(x;\theta))\thinspace\big].
\end{equation}
\end{theorem}

\begin{proof}
Placeholder proof. Show every intermediate step; cite \cref{eq:main} by
reference.
\end{proof}

\subsection{Second subsection}\label{sec:second-sub}

Placeholder derivation leading to the section's final result, which is framed:
\[
  \keyresult{\theta^\ast \;=\; \arg\min_\theta F(\theta).}
\]

\begin{corollary}[name of the corollary]\label{cor:first}
Placeholder consequence that follows from \cref{thm:main} in a line or two.
\end{corollary}

\begin{figure}
\centering
\fbox{\parbox[c][5.5cm][c]{0.96\linewidth}{\centering\color{black!45}
  \texttt{\textbackslash includegraphics\{figures/name.png\}}}}
\caption{Placeholder caption, which is the figure's interpretation. Its first
sentence names the subject or states the finding; then say what the axes
are, what each panel shows (\textbf{(a)}, \textbf{(b)}), and what pattern to
look for. \textbf{The takeaway:} the single conclusion the reader should draw.
\figcredit{scripts/plot_name.py}}
\label{fig:first}
\end{figure}

The results combine into a procedure (\cref{fig:first} illustrates it).

\begin{algorithm}[name of the procedure]\label{alg:first}
\begin{enumerate}
  \item Placeholder first step.
  \item Placeholder second step, citing the result it relies on
        (\cref{thm:main}).
  \item Placeholder final step.
\end{enumerate}
\end{algorithm}

% ===========================================================================
\part{Practice}
% ===========================================================================

\section{Case study}\label{sec:case}

Placeholder worked case: the method applied to a concrete instance, with the
measured numbers and what they mean. \textbf{State the lesson once, in bold,
and generally.}

\section{Assumptions and where they fail}\label{sec:limits}

Placeholder list of the assumptions the derivation made, and for each the
circumstances in which it breaks and what then goes wrong.

% ===========================================================================
\appendix
\section{Numerical verification}\label{app:checks}

Every analytic claim above was checked numerically; seeds are fixed so the
checks reproduce.

\begin{verification}
Placeholder claim & \ref{sec:second} & how it was checked & measured agreement \\ \rowrule
Placeholder claim & \ref{sec:second-sub} & how it was checked & measured agreement \\
\end{verification}

% ===========================================================================
\begin{thebibliography}{9}

\bibitem{ref-article} Surname, A. (Year). \textit{Title of the article.}
Journal, Volume(Issue), pages. \href{https://doi.org/}{Link}

\bibitem{ref-book} Surname, A., \& Surname, B. (Year). \textit{Title of the
Book}, edition, chapter. Publisher. \href{https://doi.org/}{Link}

\end{thebibliography}

\end{document}
